% CV, Spanish, backend. Build and public (no phone) version: see README.md in this directory.
\documentclass[10pt,a4paper]{article}
\usepackage{iftex}
\ifPDFTeX
  \usepackage[T1]{fontenc}
  \usepackage[utf8]{inputenc}
  \usepackage{lmodern}
  \input{glyphtounicode}
  \pdfgentounicode=1
\else
  \usepackage{fontspec}
\fi
\usepackage[spanish,es-noshorthands]{babel}
\usepackage[a4paper,margin=1.3cm,top=1.0cm,bottom=0.8cm]{geometry}
\usepackage{enumitem}
\usepackage{titlesec}
\usepackage{xurl}
\usepackage[hidelinks,pdfusetitle]{hyperref}
\newif\ifCVPublic
\ifdefined\CVPublic \CVPublictrue \else \CVPublicfalse \fi
\pagestyle{empty}
\setlength{\parindent}{0pt}
\setlength{\parskip}{0pt}
\setlength{\emergencystretch}{2em}
\urlstyle{same}
\raggedright
\titleformat{\section}{\large\bfseries}{}{0pt}{}[\vspace{1pt}\titlerule]
\titlespacing*{\section}{0pt}{4.5pt}{2.5pt}
\setlist[itemize]{leftmargin=1.1em,itemsep=0.6pt,topsep=1.5pt,parsep=0pt,partopsep=0pt,label=\textbullet}
% Dates stay inline after the role (no \hfill), so plain-text extraction keeps them on the role's line.
% \role{Title}{Dates}
\newcommand{\role}[2]{\textbf{#1} · #2\par}
% \proj{Name}{what it is and stack}{links}
\newcommand{\proj}[3]{\textbf{#1} · #2 · #3\par}
\newcommand{\gh}[1]{\href{https://github.com/Mateocas1/#1}{github.com/Mateocas1/#1}}
\hypersetup{
  pdftitle={Mateo Vazquez Castiglione - CV - Desarrollador Backend, TypeScript, Node.js, PostgreSQL, Integraciones con IA},
  pdfauthor={Mateo Vazquez Castiglione},
  pdfkeywords={Backend, TypeScript, Node.js, PostgreSQL, OpenAPI, LLM, MCP, evals, Python, FastAPI, Go, Java, Spring Boot, Vitest, GitHub Actions, Docker}
}
\begin{document}
{\LARGE\bfseries Mateo Vazquez Castiglione}\par
\vspace{2pt}
{\bfseries Desarrollador Backend · TypeScript · Node.js · PostgreSQL · Integraciones con IA}\par
\vspace{2pt}
Boulogne, Buenos Aires, Argentina · \ifCVPublic\else +54 9 11 4888-4832 · \fi
\href{mailto:castiglionemateo1@gmail.com}{castiglionemateo1@gmail.com}\par
\href{https://www.linkedin.com/in/vazquezmateo}{linkedin.com/in/vazquezmateo} ·
\href{https://github.com/Mateocas1}{github.com/Mateocas1} ·
\href{https://mateocas1.github.io/portfolio/}{mateocas1.github.io/portfolio}\par

\section{Resumen}
Desarrollador backend e integraciones con IA: TypeScript/Node y PostgreSQL, también Python, Go y Java. En ofertaSUPER bajé la CPU por página de 10-16\,s a 18-173\,ms (réplica local) con un índice por EAN.

\section{Experiencia}
\role{Desarrollador freelance}{may 2026 -- actualidad}
\textbf{Mena Nails}, estudio de uñas en Boulogne · remoto · may 2026 -- oct 2026 · trabajo terminado, en pre-lanzamiento. Código público: \gh{receipt-verifier} · \gh{mena-assistant}\par
\begin{itemize}
  \item Construí la API de inventario en Go y PostgreSQL (pgx, sqlc) con contrato OpenAPI 3.1, 17 migraciones, rol de DB sin ownership y tokens de servicio guardados como hash SHA-256.
  \item Diseñé el asistente de Telegram en TypeScript (grammY, MCP SDK) para que el LLM nunca escribiera solo: cada escritura era una propuesta que la dueña confirmaba con un botón de un solo uso.
  \item Elegí el modelo del asistente con un eval propio de 87 casos de tool-calling sobre 6 LLMs: el elegido aprobó completos 84 de 87.
  \item Hice un verificador de comprobantes en Python (FastAPI, Pydantic) donde aprueba código determinístico, nunca el modelo: 0 de 30 adversariales aprobados en un set sintético de 150, con los modelos elegidos.
\end{itemize}

\section{Proyectos}
\proj{ofertaSUPER}{compara precios de 5 supermercados sobre unos 24.500 productos · TypeScript, Next.js, PostgreSQL}{\gh{ofertaSUPER}}
\begin{itemize}
  \item Vercel pausó el sitio por consumo de CPU. Encontré la causa (un \texttt{filter} por producto) y la resolví con un índice por EAN: de 10-16\,s a 18-173\,ms de CPU por página, medido en réplica local.
  \item Antes de cambiar el algoritmo dejé el viejo como oráculo: salida idéntica en 25 de 25 casos sobre el snapshot real.
  \item Pipeline diario en GitHub Actions que concilia 5 supermercados en PostgreSQL bajo advisory lock, con identidad GTIN-14. Verde todos los días desde el 27/10.
\end{itemize}
\proj{Rinde tu Plata}{dónde rinden más los pesos frente a la inflación esperada · Java 21, Spring Boot 4.1}{\gh{rindetuplata}}
\begin{itemize}
  \item Ingesta con Spring Batch desde 9 fuentes: re-ejecutar un día no duplica datos (un job instance por fecha) y una fuente caída no frena a las demás.
  \item Outbox hacia RabbitMQ con confirms, consumidor idempotente y DLQ (at-least-once con deduplicación), probado con 905 tests backend contra PostgreSQL y RabbitMQ reales.
\end{itemize}

\section{Habilidades}
\textbf{Lenguajes:} TypeScript, SQL, Python, Go, Java 21\par
\textbf{Backend y datos:} Node.js, REST, OpenAPI 3.1, Zod, PostgreSQL (migraciones, advisory locks), Prisma, FastAPI, Go (pgx, sqlc), Spring Boot, RabbitMQ\par
\textbf{IA:} LLMs con tool-calling y MCP, confirmación humana, evals propios\par
\textbf{Testing y CI:} Vitest, node:test, pytest, JUnit 5, Testcontainers, GitHub Actions, Docker\par
\textbf{Cómo trabajo:} diseño, decido y verifico: defino los tests, los criterios de revisión y las mediciones. Agentes de IA escriben buena parte del código bajo esos tests y esa revisión.\par
\textbf{Frontend:} React, Next.js, Astro, Svelte\par

\section{Otra experiencia}
\role{Carrefour (Olivos) · Repositor de salón, part-time}{feb 2026 -- actualidad}
\role{Adecco (en Carrefour, 24 h semanales) · Picking con RF/PDA, 99\,\% de exactitud}{sep 2025 -- feb 2026}
\role{Mostaza · Empleado experto: caja, apertura, cierre y arqueo}{jun 2023 -- ene 2025}

\section{Formación}
\role{CBC de Ingeniería en Informática · UBA XXI}{ago 2026 -- 1.er cuatrimestre 2027; después, Ingeniería en Informática (UBA)}
\role{Diplomatura Full-Stack Developer · UTN FRBA (187 h, finalizada)}{nov 2025 -- may 2026}
\role{Secundario completo · EES N°14 Benito Lynch}{2018 -- 2023}

\section{Idiomas}
Español: nativo · Inglés: C1

\end{document}
